\section{言葉から本文へ戻る}\label{節：用語の対応}
本文で見た操作や性質には、一般に使われる名前がある。
ここには、その意味を短く記した。詳しい定義・条件・計算は、それぞれの参照先にある。
上から順に覚えるための一覧ではなく、必要な言葉から本文へ戻るための対応表である。
\begin{description}[style=nextline,leftmargin=0pt,labelsep=0pt,itemsep=5pt]
\item[\hypertarget{term:recurrence}{漸化式}] 項どうしの関係によって数列を定める式。 \hyperref[節：数列と漸化式]{本文の説明へ。}
\item[\hypertarget{term:sequence}{数列}] 数を順に並べたもの。添字で各項の位置を表す。 \hyperref[節：数列と漸化式]{本文の説明へ。}
\item[\hypertarget{term:initial}{初期条件}] 数列を作り始めるために与える値。 \hyperref[節：数列と漸化式]{本文の説明へ。}
\item[\hypertarget{term:general}{一般項}] 番号からその位置の項を直接表す式。 \hyperref[節：数列と漸化式]{本文の説明へ。}
\item[\hypertarget{term:sigma}{シグマ}] 和を、添字の範囲とともに書く記号。 \hyperref[節：シグマの計算]{本文の説明へ。}
\item[\hypertarget{term:product}{総積}] 指定した範囲の項をすべて掛けたもの。 \hyperref[節：階比型]{本文の説明へ。}
\item[\hypertarget{term:partial}{部分分数分解}] 一つの分数を、分母がより簡単な分数の和へ直す操作。 \hyperref[節：部分分数分解]{本文の説明へ。}
\item[\hypertarget{term:square}{平方完成}] 同じ数を足して引くことで、二乗の形を作る操作。 \hyperref[節：+a-aをする]{本文の説明へ。}
\item[\hypertarget{term:arithmetic}{等差数列}] 次の項へ進むとき、同じ数を足す数列。 \hyperref[節：等差型]{本文の説明へ。}
\item[\hypertarget{term:geometric}{等比数列}] 次の項へ進むとき、同じ数を掛ける数列。 \hyperref[節：等比型]{本文の説明へ。}
\item[\hypertarget{term:difference}{階差}] 隣り合う項の差。変化を足し戻すと元の数列に戻る。 \hyperref[節：階差型]{本文の説明へ。}
\item[\hypertarget{term:ratio}{階比}] 隣り合う項の倍率を追う見方。積による表示と正規化を扱う。 \hyperref[節：階比型]{本文の説明へ。}
\item[\hypertarget{term:auxiliary}{補助数列}] 元の数列から、変化を調べやすい別の量を作った数列。 \hyperref[節：補助数列の設計]{本文の説明へ。}
\item[\hypertarget{term:normalize}{正規化}] 基準となる量で割るなどして、倍率や大きさをそろえる操作。 \hyperref[節：補助数列の設計]{本文の説明へ。}
\item[\hypertarget{term:fixed}{不動点}] 規則を適用しても変わらない値。 \hyperref[節：特性方程式型]{本文の説明へ。}
\item[\hypertarget{term:characteristic}{特性方程式}] 一定の値や倍率を使って、元の式を扱いやすくするために解く方程式。本書では不動点を求める式にもこの名称を用いる。 \hyperref[節：特性方程式型]{本文の説明へ。}
\item[\hypertarget{term:root}{特性根}] 特性方程式の解。 \hyperref[節：二次特性方程式型]{本文の説明へ。}
\item[\hypertarget{term:undetermined}{未定係数法}] 式の形を先に選び、係数を比較して未知の係数を決める方法。 \hyperref[節：階差等比複合型]{本文の説明へ。}
\item[\hypertarget{term:iterate}{反復}] 同じ規則を繰り返し適用すること。 \hyperref[節：反復・不動点]{本文の説明へ。}
\item[\hypertarget{term:induction}{数学的帰納法}] 出発点と、次の番号へ進めることを示して、すべての番号について証明する方法。 \hyperref[節：数学的帰納法]{本文の説明へ。}
\item[\hypertarget{term:linear}{線形}] 入力の和と定数倍を保つ性質。線形漸化式という名称は、未知の項が係数を掛けて足す形で現れる方程式に使う。 \hyperref[節：線形・非線形]{本文の説明へ。}
\item[\hypertarget{term:nonlinear}{非線形}] 線形の条件を満たさないもの。項の二乗や逆数を含む例を扱う。 \hyperref[節：線形・非線形]{本文の説明へ。}
\item[\hypertarget{term:affine}{アフィン}] 線形な規則に定数を加えた形。 \hyperref[節：アフィン・同次・非同次]{本文の説明へ。}
\item[\hypertarget{term:homogeneous}{同次}] 線形な式で、未知の項を集めた後に付加項が残らない形。 \hyperref[節：アフィン・同次・非同次]{本文の説明へ。}
\item[\hypertarget{term:inhomogeneous}{非同次}] 線形な式で、未知の項を集めた後に既知の付加項が残る形。 \hyperref[節：アフィン・同次・非同次]{本文の説明へ。}
\item[\hypertarget{term:superposition}{重ね合わせの原理}] 同次線形の式を満たす解を、定数倍して足しても解になる性質。 \hyperref[節：線形結合]{本文の説明へ。}
\item[\hypertarget{term:combination}{線形結合}] 数列やベクトルに定数を掛けて足したもの。 \hyperref[節：線形結合]{本文の説明へ。}
\item[\hypertarget{term:independent}{一次独立}] 定数倍して足したものが零になるのは、係数がすべて零の場合だけ、という関係。 \hyperref[節：一次独立・基底・次元]{本文の説明へ。}
\item[\hypertarget{term:basis}{基底}] すべてを定数倍と足し算で表せる、一次独立な組。 \hyperref[節：一次独立・基底・次元]{本文の説明へ。}
\item[\hypertarget{term:dimension}{次元}] 基底に必要な対象の個数。解全体では、自由に選べる初期値の個数とつながる。 \hyperref[節：一次独立・基底・次元]{本文の説明へ。}
\item[\hypertarget{term:vector-space}{ベクトル空間}] 足し算と定数倍ができ、これらの操作が所定の規則を満たす集合。具体的な解の集合から説明する。 \hyperref[節：一次独立・基底・次元]{本文の説明へ。}
\item[\hypertarget{term:constant}{定数係数}] 未知の項に掛かる係数が、添字によらず一定であること。 \hyperref[節：定数係数・変数係数]{本文の説明へ。}
\item[\hypertarget{term:variable}{変数係数}] 未知の項に掛かる係数が、添字によって変わること。 \hyperref[節：定数係数・変数係数]{本文の説明へ。}
\item[\hypertarget{term:order}{階数}] 次の項を求める際に、どこまで前の項を使うかを表す数。適用する形と範囲は参照先で確認する。 \hyperref[節：階数と初期条件]{本文の説明へ。}
\item[\hypertarget{term:elementary}{初等関数}] 多項式・指数・対数・三角関数などを、有限回の演算や合成で組み合わせた関数。 \hyperref[節：初等関数・閉じた形]{本文の説明へ。}
\item[\hypertarget{term:closed}{閉じた形}] 番号が増えても式中の演算の個数を増やさずに値を表す言い方。使える関数は文脈による。 \hyperref[節：初等関数・閉じた形]{本文の説明へ。}
\item[\hypertarget{term:logarithm}{対数}] ある数が底の何乗かを表す数。真数と底の条件を伴う。 \hyperref[節：対数]{本文の説明へ。}
\item[\hypertarget{term:inverse}{逆関数}] 入力と出力の対応を逆にたどる関数。 \hyperref[節：対数]{本文の説明へ。}
\item[\hypertarget{term:binomial}{二項定理}] 二つの数の和の整数乗を、項の和へ展開する公式。 \hyperref[節：二項定理]{本文の説明へ。}
\item[\hypertarget{term:trigonometric}{三角関数}] 角度を、円上の位置や辺の比と結び付ける関数。 \hyperref[節：三角関数]{本文の説明へ。}
\item[\hypertarget{term:addition}{加法定理}] 角度の和を、元の角度の三角関数で表す公式。 \hyperref[節：三角関数]{本文の説明へ。}
\item[\hypertarget{term:limit}{極限}] 番号や変数を近づけるときに、値が近づく先を調べる考え方。 \hyperref[節：極限]{本文の説明へ。}
\item[\hypertarget{term:one-sided}{片側極限}] 変数を片側から近づけたときの極限。 \hyperref[節：片側極限]{本文の説明へ。}
\item[\hypertarget{term:left-limit}{左側極限}] 小さい側から変数を近づけたときの極限。 \hyperref[節：片側極限]{本文の説明へ。}
\item[\hypertarget{term:right-limit}{右側極限}] 大きい側から変数を近づけたときの極限。 \hyperref[節：片側極限]{本文の説明へ。}
\item[\hypertarget{term:polar}{極形式}] 複素数を、絶対値と角度によって表す形。 \hyperref[節：複素平面と極形式]{本文の説明へ。}
\item[\hypertarget{term:complex-plane}{複素平面}] 複素数を平面上の点として表す見方。 \hyperref[節：複素平面と極形式]{本文の説明へ。}
\item[\hypertarget{term:conjugate}{共役}] 複素数の虚部の符号を反転させる関係。 \hyperref[節：複素平面と極形式]{本文の説明へ。}
\item[\hypertarget{term:argument}{偏角}] 複素平面で、正の実軸から測る角度。角度の範囲の取り方を伴う。 \hyperref[節：複素平面と極形式]{本文の説明へ。}
\item[\hypertarget{term:de-moivre}{ド・モアブルの定理}] 複素数の整数乗で、絶対値と角度がどう変わるかを表す関係。 \hyperref[節：ド・モアブルの定理]{本文の説明へ。}
\item[\hypertarget{term:euler}{オイラーの公式}] 複素指数関数と正弦・余弦を結ぶ公式。 \hyperref[節：オイラーの公式]{本文の説明へ。}
\item[\hypertarget{term:matrix}{行列}] 数を長方形に並べ、複数の量の更新をまとめて表す道具。 \hyperref[節：行列の基本]{本文の説明へ。}
\item[\hypertarget{term:column-vector}{列ベクトル}] 複数の数を縦に並べた、一組の記録。 \hyperref[節：行列の基本]{本文の説明へ。}
\item[\hypertarget{term:row-vector}{行ベクトル}] 複数の数を横に並べたもの。 \hyperref[節：行列の基本]{本文の説明へ。}
\item[\hypertarget{term:identity}{単位行列}] 掛けても元のベクトルや行列を変えない行列。 \hyperref[節：行列の基本]{本文の説明へ。}
\item[\hypertarget{term:eigenvector}{固有ベクトル}] 行列を掛けると、元のベクトルの定数倍になる非零のベクトル。 \hyperref[節：行列の基本]{本文の説明へ。}
\item[\hypertarget{term:left-eigen}{左固有ベクトル}] 行列を右から掛けると、自身の定数倍になる非零の行ベクトル。 \hyperref[節：行列の基本]{本文の説明へ。}
\item[\hypertarget{term:eigenvalue}{固有値}] 固有ベクトルが更新されるときの倍率。 \hyperref[節：行列の基本]{本文の説明へ。}
\item[\hypertarget{term:transpose}{転置}] 行と列を入れ替える操作。 \hyperref[節：行列の基本]{本文の説明へ。}
\item[\hypertarget{term:invariant}{不変量}] 更新しても値が変わらない量。 \hyperref[節：不変量から線形の関係]{本文の説明へ。}
\item[\hypertarget{term:double}{二重数列}] 二つの添字で項を指定する数列。 \hyperref[節：二重数列と二方向の階差]{本文の説明へ。}
\item[\hypertarget{term:periodic}{周期性}] 一定の間隔を進めると同じ状態に戻る性質。 \hyperref[節：非線形の周期性]{本文の説明へ。}
\item[\hypertarget{term:floor}{ガウス記号}] ある数を超えない最大の整数を表す記号。 \hyperref[節：ガウス型]{本文の説明へ。}
\item[\hypertarget{term:floor-function}{床関数}] ある数を超えない最大の整数を返す関数。ガウス記号と結び付く。 \hyperref[節：ガウス型]{本文の説明へ。}
\item[\hypertarget{term:newton}{ニュートン法}] 接線の交点から、方程式の解に近い次の値を作る方法。 \hyperref[節：ニュートン法型]{本文の説明へ。}
\item[\hypertarget{term:pell}{ペル方程式}] 整数解を求める二次方程式の族。本書では二乗の差が1となる例を追う。 \hyperref[節：ペル方程式と近似分数]{本文の説明へ。}
\item[\hypertarget{term:continued}{連分数}] 分母の中に分数が続く表示。近似分数と漸化式を結び付ける。 \hyperref[節：ペル方程式と近似分数]{本文の説明へ。}
\item[\hypertarget{term:difference-operator}{差分}] 添字を一つずらした値から、元の値を引く操作。 \hyperref[節：差分から始まる離散的な微積分]{本文の説明へ。}
\item[\hypertarget{term:generating}{母関数}] 数列の項を係数に並べ、数列全体を一つの式として扱う道具。 \hyperref[節：数列全体を一つの対象にまとめると,何が見えるか]{本文の説明へ。}
\item[\hypertarget{term:convolution}{畳み込み}] 添字の和を合わせた項どうしの積を足す操作。母関数の積に現れる。 \hyperref[節：演習の操作を別の表現で見直す]{本文の説明へ。}
\item[\hypertarget{term:chebyshev}{チェビシェフ多項式}] 余弦の整数倍の公式を、多項式の列として表したもの。 \hyperref[節：特性方程式を行列で見直す]{本文の説明へ。}
\item[\hypertarget{term:legendre}{ルジャンドル多項式}] 三項漸化式で作られる多項式列。本文では直交性も紹介する。 \hyperref[節：特性方程式を行列で見直す]{本文の説明へ。}
\item[\hypertarget{term:interpolation}{ラグランジュ補間}] 有限個の点を通る多項式を作る方法。 \hyperref[節：特性方程式を行列で見直す]{本文の説明へ。}
\end{description}
